\documentclass[11pt,a4paper]{article}
\usepackage[T1]{fontenc}
\usepackage[utf8]{inputenc}
\usepackage{lmodern}
\usepackage[indonesian]{babel}
\usepackage[a4paper,margin=2.2cm]{geometry}
\usepackage{microtype}
\usepackage{amsmath,amssymb}
\usepackage{booktabs,tabularx,array}
\usepackage{enumitem}
\usepackage[most]{tcolorbox}
\usepackage{xcolor}
\usepackage{fancyhdr}
\usepackage{needspace}
\usepackage[hidelinks]{hyperref}
\setlength{\parindent}{0pt}
\setlength{\parskip}{5pt}
\setlist{itemsep=2pt,topsep=3pt}
\renewcommand{\arraystretch}{1.25}
\definecolor{biru}{HTML}{1F4E79}
\definecolor{hijau}{HTML}{2E7D32}
\definecolor{merah}{HTML}{B71C1C}
\definecolor{kuning}{HTML}{8D6E00}
\newtcolorbox{intinya}{colback=biru!6,colframe=biru,title=Intinya,fonttitle=\bfseries,boxrule=0.6pt,arc=2pt}
\newtcolorbox{hafal}{colback=hijau!6,colframe=hijau,title=Wajib hafal,fonttitle=\bfseries,boxrule=0.6pt,arc=2pt}
\newtcolorbox{awas}{colback=merah!5,colframe=merah,title=Jangan salah ucap,fonttitle=\bfseries,boxrule=0.6pt,arc=2pt}
\newtcolorbox{contoh}[1][]{colback=kuning!5,colframe=kuning,title={Contoh hitung #1},fonttitle=\bfseries,boxrule=0.6pt,arc=2pt}
\newcounter{tanya}
\newenvironment{tanya}[1]{\refstepcounter{tanya}\par\medskip\needspace{4\baselineskip}\noindent\textbf{\color{biru}T\thetanya.\ #1}\par\noindent\textbf{Jawab:}\ }{\par}
\pagestyle{fancy}\fancyhf{}
\fancyhead[L]{\small Bahan Belajar Proposal --- Abd. Rouf (223140118)}
\fancyhead[R]{\small\thepage}
\renewcommand{\headrulewidth}{0.4pt}
\newcommand{\judul}[2]{\begin{center}{\LARGE\bfseries\color{biru}#1}\par\vspace{4pt}{\large #2}\par\vspace{2pt}{\small Pegangan pribadi untuk bimbingan dan ujian proposal. Isi mengikuti \textit{Proposal\_Skripsi\_ABD\_ROUF.docx}.}\end{center}\vspace{-4pt}\hrule\vspace{6pt}}

\begin{document}
\judul{Bab II --- Tinjauan Pustaka}{Teori yang dipakai, rumusnya, dan cara menjelaskannya}

\begin{intinya}
Bab II menjawab empat pertanyaan secara berurutan:
\textbf{(1) Apa objeknya?} aplikasi presensi \textit{offline-first} (2.1.1).
\textbf{(2) Apa yang diukur?} QoS, TIPHON, jaringan akses seluler \& Wi-Fi (2.1.2--2.1.4).
\textbf{(3) Bagaimana mengukurnya?} \textit{crowdsourcing} dan telemetri (2.1.5--2.1.6).
\textbf{(4) Bagaimana menganalisisnya?} K-Means, penentuan $k$, Kruskal-Wallis (2.1.7--2.1.9),
ditambah metode penelitian (2.1.10) dan pelindungan data pribadi (2.1.11). Lalu 11 penelitian terdahulu (2.2).
\end{intinya}

\section*{1. Objek: aplikasi presensi dan \textit{offline-first} (2.1.1)}
\begin{itemize}
  \item \textbf{Offline-first} = data disimpan dulu di HP (Isar) sebagai sumber utama, dimasukkan ke antrean, lalu dikirim ke server (PocketBase via HTTPS) saat ada jaringan (Devina dkk., 2026).
  \item Penting untuk penelitian: \textbf{antrean mencatat setiap percobaan kirim dan hasilnya}, jadi bisa \textit{diamati} tanpa diubah.
  \item \textbf{Geofence}: presensi hanya bisa dalam radius 15--30 m lokasi kerja. Pengecualian: GAMAS (tugas luar), lokasi fleksibel dari admin, titik sementara admin. Geofence tetap jalan tanpa internet (GPS + daftar lokasi tersimpan di HP).
  \item Wi-Fi kantor \textbf{tidak wajib}; aplikasi hanya mencatat apakah Wi-Fi yang tersambung itu Wi-Fi kantor (berdasarkan BSSID).
  \item Kenapa dibahas? Karena geofence menentukan \textbf{apakah lokasi suatu pengukuran bisa dipercaya} (dipakai saat membentuk sesi di Bab III).
\end{itemize}

\section*{2. Parameter QoS (2.1.2)}
\begin{center}
\begin{tabularx}{\linewidth}{@{}p{2.2cm}Xp{4.6cm}@{}}
\toprule
Parameter & Arti sederhana & Rumus \\
\midrule
Delay (RTT) & Waktu kirim permintaan sampai balasan diterima (pulang-pergi) & $d = t_{terima} - t_{kirim}$ \hfill (2.1)\\
Jitter & Seberapa ``goyang'' delay dari satu pengukuran ke berikutnya & $J=\frac{1}{n-1}\sum_{i=2}^{n}|d_i-d_{i-1}|$ \hfill (2.2)\\
Packet loss & Persen paket yang tidak sampai & $PL=\frac{N_{kirim}-N_{terima}}{N_{kirim}}\times100\%$ \hfill (2.3)\\
Throughput & Kecepatan nyata kirim data & $T=\frac{B\times 8}{\Delta t}$ \hfill (2.4)\\
\bottomrule
\end{tabularx}
\end{center}
Satuan throughput: $B$ dalam byte, $\Delta t$ dalam ms, hasil dalam \textbf{kbps} (karena bit per milidetik = kilobit per detik).

\begin{contoh}[delay, jitter, throughput]
Tiga probe ke server: RTT = 100, 120, 110 ms.\\
\textbf{Delay sesi} = median = \textbf{110 ms}.\\
\textbf{Jitter} = $\frac{|120-100|+|110-120|}{3-1}=\frac{20+10}{2}=$ \textbf{15 ms}.\\
\textbf{Throughput}: foto presensi 150.000 byte terkirim dalam 1.200 ms $\Rightarrow T=\frac{150.000\times8}{1.200}=$ \textbf{1.000 kbps} (kategori TIPHON ``sedang'', indeks 2).
\end{contoh}

\textbf{Mengapa packet loss diganti ``tingkat kegagalan''?} Aplikasi Android biasa tidak bisa melihat paket tanpa akses khusus (root/VPN). Maka dihitung di tingkat aplikasi: persen \textit{permintaan} yang gagal karena \textit{timeout} atau galat jaringan. Sama-sama menggambarkan ``data tidak sampai'', tetapi lebih kasar karena satu permintaan = banyak paket. \textbf{Ini keterbatasan yang diakui.}

\section*{3. Standar TIPHON (2.1.3)}
\begin{hafal}
\small\begin{tabularx}{\linewidth}{@{}l>{\raggedright}X>{\raggedright}X>{\raggedright}X>{\raggedright\arraybackslash}Xc@{}}
Kategori & Delay & Jitter & Gagal/PL & Throughput & Indeks\\
\midrule
Sangat bagus & $<150$ ms & 0 ms & $<3\%$ & $>2{,}1$ Mbps & 4\\
Bagus & 150--300 & 0--75 & 3 -- $<15\%$ & 1,2--2,1 Mbps & 3\\
Sedang & 300--450 & 75--125 & 15 -- $<25\%$ & 700 kbps--1,2 Mbps & 2\\
Jelek & $>450$ & 125--225 & $\geq 25\%$ & $<700$ kbps & 1\\
\end{tabularx}

\medskip
Indeks keseluruhan = rata-rata 4 indeks: \textbf{3,80--4,00} sangat memuaskan; \textbf{3,00--3,79} memuaskan; \textbf{2,00--2,99} kurang memuaskan; \textbf{1,00--1,99} tidak memuaskan.
\end{hafal}
\begin{itemize}
  \item Di tabel asli ETSI, nilai packet loss (0\%, 3\%, 15\%, 25\%) adalah \textbf{batas bawah}; di proposal ditafsirkan jadi rentang seperti di atas.
  \item TIPHON aslinya untuk \textbf{suara/multimedia real-time}, jadi cukup ketat untuk presensi. Tetap dipakai agar \textbf{bisa dibandingkan} dengan penelitian QoS lain; pola sebenarnya juga dilihat dari K-Means yang tidak bergantung ambang TIPHON.
\end{itemize}

\section*{4. Jaringan akses seluler dan Wi-Fi (2.1.4)}
\begin{itemize}
  \item \textbf{Wi-Fi}: HP $\rightarrow$ \textit{access point} (AP) $\rightarrow$ ISP kantor $\rightarrow$ internet.
  \item \textbf{Seluler}: HP $\rightarrow$ BTS/eNodeB operator $\rightarrow$ jaringan inti operator $\rightarrow$ internet.
  \item Kualitas dipengaruhi: \textbf{jarak \& hambatan} ke AP/BTS, \textbf{jumlah pengguna} yang berbagi kapasitas, dan \textbf{kapasitas jalur} operator/ISP ke server.
  \item Akibatnya: dua orang di lokasi sama bisa beda kualitas jika operatornya beda; kualitas turun saat ramai (jam sibuk). Jadi \textbf{jenis jaringan dan operator wajib dicatat}.
  \item Nama operator SIM data dibaca dengan \texttt{TelephonyManager} (tanpa izin tambahan).
  \item Kalimat kunci: \textit{penelitian ini mengukur infrastruktur fisik jaringan (BTS dan AP) dari sisi pengguna, dengan ponsel karyawan sebagai sensor.}
\end{itemize}

\section*{5. Crowdsourcing dan telemetri (2.1.5--2.1.6)}
\begin{itemize}
  \item \textbf{Crowdsourced network measurement}: diukur di HP pengguna akhir, banyak perangkat, lokasi, dan waktu (Wamser dkk., 2021).
  \item \textbf{Participatory} = pengguna sengaja menjalankan tes. \textbf{Opportunistic} = berjalan otomatis saat ada kesempatan (Wu dkk., 2017). \textbf{Penelitian ini: opportunistic, dipicu presensi.}
  \item Kelemahan: lokasi dan waktu tidak dikendalikan. Penanganan: diringkas \textbf{per sesi}, minimal \textbf{30 sesi} per lokasi untuk uji beda, disajikan \textbf{agregat} (Wamser 2021; Tripkovic 2021).
  \item \textbf{Telemetri} = modul di satu HP yang mengukur lalu mengunggah. \textbf{Crowdsourcing} = gabungan data dari banyak HP.
  \item \textbf{Kenapa bukan ping?} Aplikasi biasa tidak bisa kirim ICMP tanpa root. Maka diukur \textbf{RTT tingkat aplikasi} (HTTP) ke \texttt{/api/health} yang sangat ringan (tidak membuka basis data), sehingga RTT didominasi waktu jaringan.
  \item \textbf{Hibrida}: pasif (sinkronisasi, unggah foto, status koneksi) + aktif (\textit{probe} kecil saat presensi).
\end{itemize}
\begin{hafal}
\textbf{Tabel 2.3 --- dua alamat pengukuran}\par
\begin{tabularx}{\linewidth}{@{}p{2.6cm}p{3.2cm}X@{}}
Server perusahaan & Pembanding (\texttt{generate\_204}) & Artinya\\
\midrule
Baik & Baik & Jaringan dan server baik\\
Buruk & Buruk & Masalah di jaringan akses/operator (BTS, AP, ISP)\\
Buruk & Baik & Masalah di server perusahaan atau jalur ke server\\
\end{tabularx}
\end{hafal}

\section*{6. K-Means (2.1.7)}
K-Means membagi $n$ data ke $k$ klaster; tiap data masuk ke pusat (\textit{centroid}) terdekat. \textit{Unsupervised} (tidak butuh label). Kelebihan: sederhana, cepat, mudah ditafsirkan. Kelemahan: $k$ harus ditentukan dan hasil dipengaruhi pusat awal.

\textbf{Langkah:} (1) tentukan $k$; (2) pilih pusat awal dengan \textbf{k-means++} (saling berjauhan, lebih stabil); (3) hitung jarak Euclidean (2.5); (4) masukkan ke pusat terdekat; (5) hitung ulang pusat = rata-rata anggota (2.6); (6) ulangi 3--5 sampai anggota tidak berubah.
\[
D(x_i,c_j)=\sqrt{\sum_{p=1}^{m}(x_{ip}-c_{jp})^2}\quad(2.5)\qquad
c_j=\frac{1}{|C_j|}\sum_{x_i\in C_j}x_i\quad(2.6)
\]
\textbf{Normalisasi wajib}, karena delay ratusan ms sedangkan kegagalan 0--100\%; tanpa normalisasi delay akan mendominasi. Dipakai \textbf{RobustScaler} (median dan IQR, tahan pencilan):
\[x'=\frac{x-Q_2}{Q_3-Q_1}\quad(2.7)\]

\section*{7. Menentukan jumlah klaster (2.1.8)}
\begin{itemize}
  \item \textbf{Elbow}: hitung SSE untuk beberapa $k$, pilih ``siku'' (setelah itu SSE tidak turun berarti). $SSE=\sum_{j}\sum_{x_i\in C_j}D(x_i,c_j)^2$ (2.8). Perlu dicek dengan pemahaman bidang (Mihevc \& Podržaj, 2025).
  \item \textbf{Silhouette}: $s(i)=\frac{b(i)-a(i)}{\max\{a(i),b(i)\}}$ (2.9). $a(i)$ = rata-rata jarak ke teman sekelompok; $b(i)$ = rata-rata jarak ke kelompok lain terdekat. Rentang $-1$ s.d. $1$: dekat 1 = tepat, $\approx 0$ = di perbatasan, negatif = mungkin salah klaster. Pembanding: Irwansyah (2025) memperoleh 0,351.
\end{itemize}
\begin{contoh}[silhouette]
Sebuah sesi punya $a(i)=0{,}3$ dan $b(i)=1{,}2$. Maka $s(i)=\frac{1{,}2-0{,}3}{1{,}2}=0{,}75$ $\Rightarrow$ sesi itu jelas berada di klaster yang tepat.
\end{contoh}

\section*{8. Uji Kruskal-Wallis (2.1.9)}
\begin{itemize}
  \item Uji \textbf{nonparametrik}: apakah 3+ kelompok independen punya median sama. Untuk 2 kelompok setara \textbf{Mann-Whitney}.
  \item Tidak butuh data normal karena memakai \textbf{peringkat}. Cocok untuk delay yang \textbf{menceng ke kanan} (banyak kecil, sedikit sangat besar).
  \item $H=\frac{12}{N(N+1)}\sum_{j=1}^{g}\frac{R_j^2}{n_j}-3(N+1)$ (2.10), dibandingkan dengan \textbf{chi-square}, derajat bebas $g-1$. Jika $p<0{,}05$ $\Rightarrow$ minimal satu kelompok berbeda.
  \item Lanjut \textbf{uji Dunn + koreksi Bonferroni} untuk mencari pasangan mana yang beda (koreksi diperlukan karena banyak perbandingan menaikkan peluang ``beda palsu'').
\end{itemize}
\begin{contoh}[Kruskal-Wallis, data latihan]
Delay (ms) tiga lokasi: A = 80, 95, 120; B = 150, 210, 300; C = 100, 130, 400. $N=9$, $g=3$.\\
Urutkan semua dan beri peringkat: 80(1) 95(2) 100(3) 120(4) 130(5) 150(6) 210(7) 300(8) 400(9).\\
$R_A=1+2+4=7$, \ $R_B=6+7+8=21$, \ $R_C=3+5+9=17$ (cek: total 45 = $9\cdot10/2$).\\
$H=\frac{12}{9\cdot10}\left(\frac{7^2}{3}+\frac{21^2}{3}+\frac{17^2}{3}\right)-3\cdot10=0{,}1333\times259{,}67-30=\mathbf{4{,}62}$.\\
Nilai kritis chi-square ($db=2$, $\alpha=0{,}05$) = 5,991. Karena $4{,}62<5{,}991$ ($p\approx0{,}099$), \textbf{H$_0$ diterima}: belum terbukti berbeda. Pelajaran: dengan data sedikit, selisih yang kelihatan besar pun bisa tidak signifikan --- itulah alasan minimal 30 sesi.
\end{contoh}

\section*{9. Metode penelitian dan data pribadi (2.1.10--2.1.11)}
\begin{itemize}
  \item \textbf{Kuantitatif}: data angka dari instrumen, dianalisis statistik (Sugiyono, 2019).
  \item \textbf{Deskriptif}: menggambarkan nilai satu variabel. \textbf{Komparatif}: membandingkan variabel pada dua atau lebih sampel.
  \item Variabel = kualitas jaringan pada sesi presensi; dibandingkan antarlokasi \& antaroperator; instrumen = modul telemetri yang divalidasi (harus \textbf{valid} dan \textbf{reliabel}).
  \item \textbf{UU PDP 27/2022}: pemrosesan harus terbatas, spesifik, sah, transparan.
  \item \textbf{Pseudonim}: 16 karakter awal SHA-256(ID pengguna + \textit{salt} rahasia). Beda dengan anonim karena \textbf{masih bisa dikaitkan kembali} oleh pemegang \textit{salt} $\Rightarrow$ tetap diperlakukan sebagai data pribadi.
\end{itemize}

\section*{10. Penelitian terdahulu (2.2) --- tiga kelompok}
\begin{center}
\begin{tabularx}{\linewidth}{@{}p{3.1cm}Xp{4.4cm}@{}}
\toprule
Kelompok & Isi & Kekurangannya (celah)\\
\midrule
1. QoS manual (no. 1--5) & Rosid 2023; Fiqria \& Purnama 2024; Ainina \& Ridwan 2024; Basmar \& Yusuf 2026; Makarim 2026 --- Wireshark/nPerf/SpeedTest + TIPHON atau uji statistik & Manual, satu titik, satu waktu\\
2. Crowdsourcing \& K-Means (no. 6--9) & Wu 2017 (MopEye); Wamser 2021; Mubarak 2023 (Cacti, 32 OPD); Irwansyah 2025 (Ookla, 17 kab/kota, silhouette 0,351) & Data dari aplikasi khusus, pihak ketiga, atau alat pemantau\\
3. Aplikasi presensi (no. 10--11) & Indah \& Wati 2025 (offline QR); Rokhman 2026 (geofence, \texttt{client\_event\_id}) & Tangani koneksi buruk tanpa mengukur jaringan\\
\bottomrule
\end{tabularx}
\end{center}
\textbf{Kalimat penutup yang harus bisa diucapkan:} ``Penelitian ini menggabungkan ketiganya: mengukur QoS otomatis dari aplikasi presensi, menilainya dengan TIPHON, menguji perbedaan dengan Kruskal-Wallis, dan mengelompokkan dengan K-Means per lokasi dan operator.''

\section*{11. Kemungkinan pertanyaan penguji}
\setcounter{tanya}{0}
\begin{tanya}{RTT tingkat aplikasi kan termasuk waktu server, berarti bukan murni jaringan?}
Benar, karena itu probe diarahkan ke \texttt{/api/health} yang tidak membuka basis data sehingga waktu server sangat kecil. Selain itu ada permintaan pembuka (TCP+TLS) sehingga yang diukur adalah RTT pada koneksi yang sudah terbuka, dan hasilnya divalidasi terhadap PCAPdroid/Wireshark.
\end{tanya}
\begin{tanya}{Kenapa pakai median, bukan rata-rata?}
Delay menceng ke kanan; satu nilai sangat besar menarik rata-rata jauh. Median lebih mewakili kondisi umum.
\end{tanya}
\begin{tanya}{Kenapa tidak ANOVA?}
ANOVA mensyaratkan data normal dan varians sama. Data delay umumnya tidak normal, jadi dipakai Kruskal-Wallis yang berbasis peringkat.
\end{tanya}
\begin{tanya}{Kenapa RobustScaler, bukan Min-Max atau Z-score?}
Min-Max dan Z-score sangat terpengaruh pencilan (satu delay 5.000 ms bisa ``memampatkan'' data lain). RobustScaler memakai median dan IQR sehingga tahan pencilan; dipakai juga oleh Irwansyah (2025) pada data internet seluler.
\end{tanya}
\begin{tanya}{Kenapa generate\_204?}
Alamat milik Google yang juga dipakai Android sendiri untuk mengecek internet; sangat ringan (balasan kosong) dan servernya tersebar luas, jadi cocok sebagai patokan ``internet umum''.
\end{tanya}
\begin{tanya}{TIPHON itu untuk VoIP, cocok untuk presensi?}
Ambangnya memang ketat. TIPHON dipakai sebagai tolok ukur umum agar bisa dibandingkan dengan penelitian lain; keputusan pola tidak hanya dari TIPHON tetapi juga K-Means.
\end{tanya}
\begin{tanya}{Apa bedanya pseudonim dan anonim?}
Anonim tidak bisa dikembalikan ke orangnya. Pseudonim masih bisa jika punya \textit{salt}. Karena itu data tetap diperlakukan sebagai data pribadi (akses admin, salt tidak diekspor, hasil agregat).
\end{tanya}

\begin{awas}
\begin{itemize}
  \item Satuan: delay \& jitter dalam \textbf{ms}, throughput dalam \textbf{kbps/Mbps}, kegagalan dalam \textbf{\%}.
  \item Kruskal-Wallis membandingkan \textbf{median/peringkat}, bukan rata-rata.
  \item K-Means \textbf{bukan} klasifikasi (tidak ada label); label ``baik/buruk'' diberikan \textit{setelah} melihat pusat klaster.
  \item Silhouette \textbf{bukan} akurasi.
\end{itemize}
\end{awas}

\section*{12. Latihan mandiri}
\begin{enumerate}
  \item RTT probe 200, 260, 230 ms. Hitung delay (median) dan jitter. \textit{(Jawaban: 230 ms; (60+30)/2 = 45 ms.)}
  \item Sesi dengan delay 320 ms, jitter 95 ms, gagal 12,5\%, throughput 850 kbps. Berapa indeks TIPHON-nya? \textit{(Jawaban: (2+2+3+2)/4 = 2,25 = kurang memuaskan.)}
  \item Server buruk, generate\_204 baik. Di mana masalahnya? \textit{(Jawaban: di server perusahaan atau jalur ke server.)}
  \item Jelaskan dalam dua kalimat kenapa data harus dinormalisasi sebelum K-Means.
\end{enumerate}
\end{document}
